\section{Auctions on \podtm through the \bidsettm protocol}\label{sec:auctions}
In this section, we show how single-shot distributed auctions can be implemented on top of \podtm.
This is achieved through \bidsettm, a primitive for collecting a set of bids.
The idea is as follows. A pre-appointed \emph{sequencer} -- which can be any party, even a pod replica -- runs the auction, but the bids are collected from \podtm using a \bidsettm protocol. The past-perfection property of \podtm renders the sequencer unable to censor bids: when it creates an output, all timely and honestly-written bids \emph{must} be in it, otherwise the sequencer has provably misbehaved and can be held accountable.
% We first define \bidsettm and then construct it using an underlying \podtm.


\begin{definition}[\bidsettm protocol]\label{def:auction-framework}
    A \bidsettm \emph{protocol} has a \emph{starting time} parameter \tzero and exposes the following interfaces to \emph{bidder} and \emph{consumer} parties:
    \begin{itemize}
        \item function $\opSubmitBid{\bid}$: It is called by a bidder at round \tzero to submit a bid \bid.
        \item event $\opResult{\bag, \bidcert}$: It is an event generated by a consumer. It contains a \emph{bid-set} \bag, which is a set of bids, and auxiliary information \bidcert.
    \end{itemize}
\end{definition}


A \bidsettm protocol satisfies the following liveness and safety properties:
\begin{description}
    \item[(Liveness) Termination within \parAuctionTerm:]
    An honest consumer generates an event $\opResult{\bag, \bidcert}$ by round $\tzero + \parAuctionTerm$.

    \item [(Safety) Censorship resistance:]
    If an honest bidder calls \opSubmitBid{\bid}
    and an honest consumer generates an event $\opResult{\bag, $\cdot$}$,
    then $\bid \in \bag$.
    %  or a subset of the parties that run the \bidsettm protocol is accountable.
    % \hue{Bid inclusion: In a synchronous network, all honest bids are included}
    % Moreover, \emph{fair punishment} holds: an honest sequencer cannot be held accountable for censoring bids, even in asynchronous networks.
    % \oa{Notice that, even in asynchrony, an honest sequencer cannot be blamed, in the construction we present.}

    \item[(Safety) Weak consistency:]
    If two honest consumers generate \opResult{$\bag_1, \cdot$} and \opResult{$\bag_2, \cdot$} events, such that $\bag_1 \neq \emptyset$ and $\bag_2 \neq \emptyset$, then $\bag_1 = \bag_2$.
    % \hue{Single-result: If there are multiple conflicting auction results, the sequencer can be held accountable}
    % \oa{Only in synchrony.}

\end{description}

%As with previous protocols, we require that \bidsettm satisfies \emph{accountable safety}. This property means that, if a safety property is violated, then an \identify{} function returns a subset of the parties that run the \bidsettm protocol and are malicious, and that honest parties are never returned by \identify{}.

% \subsection{Construction}
%We present \bidsettm-core, a \bidsettm construction that uses a \podtm to register bids.

% We now present \bidsettm-core, a \bidsettm construction that uses \podtm-core to register bids.

\begin{protocol}[\bidsettm-core]\label{const:bidset}
	Construction \bidsettm-core is parameterized by an integer $\Delta$ (looking ahead, we will prove security in synchrony, \textit{i.e.}, assuming the network delay $\delta$ is smaller than $\Delta$) and assumes digital signatures and a \podtm with $\delta$-timeliness, $\parPerf=\delta$ and $\parConf=2\delta$. At time \tzero, bidders execute \Cref{alg:bid-set-bidder}, the sequencer \Cref{alg:bid-set-sequencer}, and the consumers \Cref{alg:bid-set-consumer}.
	% A bidder runs \opSubmitBid{} (\Cref{alg:bid-set-bidder}), the sequencer runs \opRunAuction{} (\Cref{alg:bid-set-sequencer}), and a consumer runs \opReadResult{} (\Cref{alg:bid-set-consumer}), in order to output the $\opResult{\bag, \bidcert}$ event.
	% A pre-appointed \emph{sequencer} is responsible to reading the \podtm and writing back to it when a specific condition is met. For example, when instantiating \bidsettm-core on top of \podtm-core, a replica can act as sequencer.
\end{protocol}

\vspace{-3mm}
\begin{algorithm}[]
    \caption{\bidsettm-core: Code for a bidder. It runs a client for a \podtm-core instance \podInst.}
    \label{alg:bid-set-bidder}
    \begin{algorithmic}[1]
        \Function{\opSubmitBid{\bid}}{}
            \State{\podInst.\opwrite{\bid}}
        \EndFunction
\end{algorithmic}
\end{algorithm}

\vspace{-3mm}
\begin{algorithm}[]
    \caption{\bidsettm-core: Code for the sequencer. It runs a client for a \podtm-core instance \podInst, and $\sk_a$ denotes the secret key of the sequencer.}
    \label{alg:bid-set-sequencer}
    \begin{algorithmic}[1]
        \Function{\opRunAuction{}}{}
            \Let{\left( (\txSet, \rperf), (\ppcert, \alltxcert) \right)}{\podInst.\opreadall{}}
            \While{$\rperf \leq \tzero + \Delta$}\label{line:bidset-wait-pp}
                \Let{\left( (\txSet, \rperf), (\ppcert, \alltxcert) \right)}{\podInst.\opreadall{}}
            \EndWhile{}
            \State{$\bag \gets \{\tx \mid (\tx, \cdot, \cdot , \cdot) \in \txSet\}  ;\; \bidcert \gets \ppcert$}\label{line:bidset-create-bag}
            \Let{\sigma}{\opSign{\sk_a, (\bag, \bidcert)}} \label{line:bidset-sequencer-sign}
            \Let{\tx}{\msgbidresult{(\bag, \bidcert, \sigma)}}\label{line:bidset-write-bag}
            \State{\podInst.\opwrite{\tx}} %Run by \tzero + 2 \Delta.
        \EndFunction
\end{algorithmic}
\end{algorithm}

\vspace{-3mm}
\begin{algorithm}[]
    \caption{\bidsettm-core: Code for a consumer. It runs a client for a \podtm-core instance \podInst.}
    \label{alg:bid-set-consumer}
    \begin{algorithmic}[1]
        \Function{\opReadResult{}}{}
            \Loop{}\label{line:bidset-wait-pp-consumer}
                \Let{((\txSet, \rperf), (\ppcert, \alltxcert))}{\podInst.\opreadall{}}
                    \If{$\exists (\tx, \cdot, \cdot, \rconf, \cdot) \in \txSet$ \textbf{:} $\tx = \msgbidresult{(\bag, \bidcert, \sigma)}$ \andl $\rconf \leq \tzero + 3 \Delta$} \label{line:bidset-find-confirmed-tx}
                    \State{\textbf{output event} \opResult{\bag, \bidcert}}
                \ElsIf{$\rperf > \tzero + 3 \Delta$}\label{line:bidset-read-result-pp}
                    \State{\textbf{output event} \opResult{$\emptyset$, \ppcert}}
                \EndIf
            \EndLoop\label{line:bidset-consumer-loop-end}
        \EndFunction
    \end{algorithmic}
    \end{algorithm}

% \vspace{-5mm}

A bidder (\Cref{alg:bid-set-bidder}) submits a bid by writing it on the \podtm  at round \tzero.
The sequencer (\Cref{alg:bid-set-sequencer}) waits until the \podtm returns a past-perfect round larger than $\tzero + \Delta$ (\cref{line:bidset-wait-pp}) and then constructs the bid-set \bag from the set of transactions in \txSet (\cref{line:bidset-create-bag}).
The sequencer concludes by signing \bag and \bidcert (which can be used as evidence, in case of a safety violation) and writing  \msgbidresult{(\bag, \bidcert, \sigma)} on \podInst.
The consumer (\Cref{alg:bid-set-consumer}) waits until one of the following two conditions is met.
First, a \emph{confirmed} transaction \msgbidresult{(\bag, \bidcert, \sigma)} appears in \txSet , for which $\rconf \leq \tzero + 3 \Delta$ (\cref{line:bidset-find-confirmed-tx}), in which case it outputs bid-set \bag as result.
Second, a round higher than $\tzero + 3 \Delta$ becomes past-perfect in \podtm (\cref{line:bidset-read-result-pp}) without a confirmed \msgbidresult{} transaction appearing, in which case it outputs $\bag=\emptyset$.

As an intuition on how \bidsettm-core achieves censorship resistance, we observe the following.
The \emph{$\delta$-timeliness property} of \podtm (\Cref{def:pod-timely}), given that $\delta \leq \Delta$, ensures that bids of honest parties will have a \confirmed round $\rconf \leq \tzero + \Delta$.
Now, the sequencer may only produce a valid bid-set when \podtm returns a past-perfect round larger than $\tzero + \Delta$ (\cref{line:bidset-wait-pp}), and the output \emph{must} contain a certificate \ppcert that proves this. However, if the certificate is valid, then the sequencer must have \emph{provably} seen the bids of honest parties in \txSet (we remind that the votes of replicas on \podtm-core are chained using sequence numbers), and thus \bag must contain all bids with $\rconf \leq \tzero + \Delta$. If any party presents a transaction certificate \txcert for some transaction $\tx^*$ with $\rconf^* \leq \tzero + \Delta$, but $\tx^* \not\in \bag$, then the sequencer can be held accountable.
% We show the detailed proof in \Cref{lem:bidset-censorship} and \Cref{lem:bidset-accountability}.

% Regarding liveness, \cref{line:bidset-wait-pp} of \Cref{alg:bid-set-sequencer} becomes true in the view of sequencer by round $\tzero + \Delta + \delta$ (from the \emph{past-perfection within $\parPerf=\delta$} property of \podtm-core),
% hence \Cref{alg:bid-set-sequencer} for an honest sequencer terminates by that round.
% Observe also that the transaction \msgbidresult{(\bag, \bidcert, \sigma)} becomes confirmed in the view of all honest clients by round $\tzero + \Delta + 3\delta$ (from the \emph{confirmation within $\parConf=2\delta$} property), and it will have a confirmed round $\rconf \leq \tzero + \Delta + 2\delta$ (from the \emph{$\delta$-timeliness} property).
% Hence, if the network is synchronous and the sequencer honest, the condition in \cref{line:bidset-find-confirmed-tx} of \Cref{alg:bid-set-consumer} becomes true at round at most $\tzero + \Delta + 3\delta$,
% Even if the sequencer is malicious, from the \emph{past-perfection within $\parPerf=\delta$} property of \podtm, the condition in \cref{line:bidset-read-result-pp} will become true at latest at round $\tzero + 3 \Delta + \delta$, hence \bidsettm-core achieves \emph{termination within $\parAuctionTerm=3\Delta + \delta$}.

% hence it outputs \bag as the bid-set, hence \bidsettm-core satisfies the \emph{consistency} property.




\begin{theorem}[Bidset security]\label{thm:bidset-construction-secure}
Assuming a synchronous network where $\delta \leq \Delta$,
protocol \bidsettm-core (Construction~\ref{const:bidset}) instantiated with a digital signature and a  secure \podtm protocol that satisfies the \emph{past-perfection within $\parPerf = \delta$}, \emph{confirmation within $\parConf=2\delta$} and \emph{$\delta$-timeliness} properties, is a secure \bidsettm protocol satisfying \emph{termination within $\parAuctionTerm = 3\Delta + \delta$}.
It satisfies accountable safety with an \identifySequencer{} function that identifies a malicious sequencer.
\end{theorem}
\vspace{-3mm}
\begin{proof} The proof and \identifySequencer{} are shown in Appendix~\ref{app:bidsetsec}.
\end{proof}

\begin{remark}
    Observe that \bidsettm-core terminates within $\parAuctionTerm = 3\Delta + \delta$ in the worst case, but, if the sequencer is honest, then it terminates within $\parAuctionTerm = \Delta + 3\delta$.
    Moreover, \bidsettm-core is not responsive because \Cref{alg:bid-set-sequencer} waits for a fixed $\Delta$ interval.
    This step can be optimized if the set of bidders is known (e.g., requiring them to pre-register), which allows for the protocol to be made optimistically responsive (\textit{i.e.}, $\parAuctionTerm = 4\delta$) when all parties are honest.
\end{remark}

\begin{remark}[Implicit sub-session identifiers]
	We assume that each instance of the \bidsettm-core protocol is identified by a unique sub-session identifier (ssid). All messages written to the underlying \podtm are concatenated with the ssid.
\end{remark}

\dotparagraph{\textbf{Auctions using \bidsettm}} Building on a \bidsettm protocol, it is trivial to construct single-shot first-price and second-price open auctions as follows: 1. Bidders place their open bids $\bid$ by calling \opSubmitBid{\bid}; 2. Consumers determine the winner by calling \opReadResult{} to obtain $\bag$ and outputting either the first or second highest bid. We conjecture that single-shot sealed bid auction protocols such as those of~~\cite{trustee,seal,ACNS:DavGenPou22,CCS:TAFWBM23,chitra_et_al:LIPIcs.AFT.2024.19,cryptoeprint:2024/1011} can also be instantiated on top of a \bidsettm protocol. Intuitively, this holds because such protocols first agree on a set of sealed bids and then execute extra steps to determine the winner. However, a formal analysis of sealed-bid auction protocols based on \bidsettm is left as future work.
